\documentclass[11pt,a4paper]{article}
\usepackage{brainteaser}

\begin{document}
\brainteasertitle{17}{Les 100 prisonniers et les boîtes}

\begin{statementbox}
Cent prisonniers sont numérotés de $1$ à $100$. Dans une salle se trouvent
$100$ boîtes fermées, elles aussi numérotées de $1$ à $100$. Chacune contient
le numéro d'un prisonnier, les $100$ numéros étant répartis aléatoirement, un
par boîte.

Les prisonniers entrent dans la salle un à un. Chacun peut ouvrir au plus
$50$ boîtes et doit retrouver son propre numéro. Il referme ensuite toutes les
boîtes exactement comme il les a trouvées. Toute communication est interdite
dès que l'épreuve commence, mais les prisonniers peuvent convenir auparavant
d'une stratégie commune.

Ils sont tous libérés si chacun retrouve son numéro. Si un seul échoue, ils
perdent tous.

Peuvent-ils élaborer une stratégie leur donnant plus de $30\,\%$ de chances
de réussite collective?
\end{statementbox}

\begin{hintsbox}
\begin{enumerate}
  \item Généraliser le problème à un nombre pair $n$ de prisonniers, chacun
  pouvant ouvrir $n/2$ boîtes. Étudier d'abord les cas $n=2$ et $n=4$.
  \item Pour $n=4$, un prisonnier ouvre d'abord la boîte portant son numéro.
  S'il y découvre le numéro $j$, quelle boîte peut-il ouvrir ensuite sans
  recourir au hasard?
  \item Représenter le contenu des boîtes par une permutation $\sigma$, où
  $\sigma(i)$ est le numéro placé dans la boîte $i$. Que produit l'itération
  $i,\sigma(i),\sigma^2(i),\ldots$?
  \item Relier la réussite d'un prisonnier à la longueur du cycle contenant
  son numéro.
  \item Pour $k>n/2$, compter les permutations qui possèdent un cycle de
  longueur $k$: choisir ses éléments, former le cycle, puis permuter les
  éléments restants.
  \item Deux cycles de longueurs strictement supérieures à $n/2$ peuvent-ils
  appartenir à la même permutation?
\end{enumerate}
\end{hintsbox}

\correction

Modélisons le placement par une permutation $\sigma$ de
$\{1,\ldots,100\}$: la boîte $i$ contient le numéro $\sigma(i)$.

Le prisonnier $i$ applique la stratégie suivante. Il ouvre d'abord la boîte
$i$. S'il y lit $\sigma(i)$, il ouvre ensuite la boîte $\sigma(i)$, puis la
boîte $\sigma^2(i)$, et ainsi de suite, jusqu'à trouver $i$ ou atteindre la
limite de $50$ ouvertures.

\begin{keyideabox}
Chaque prisonnier suit le cycle de la permutation qui contient son propre
numéro.
\end{keyideabox}

Si le cycle contenant $i$ a pour longueur $\ell$, les boîtes ouvertes sont
successivement
\[
i,\ \sigma(i),\ \ldots,\ \sigma^{\ell-1}(i),
\]
et la dernière contient
$\sigma^\ell(i)=i$. Le prisonnier $i$ réussit donc si et seulement si
$\ell\leq 50$. Par conséquent, les cent prisonniers réussissent tous si et
seulement si aucun cycle de $\sigma$ n'a une longueur supérieure à $50$.

\subsection*{Les petits cas}

Pour $n=2$, chaque prisonnier ouvre sa propre boîte. Ils réussissent ensemble
exactement lorsque la permutation est l'identité, soit avec probabilité
$1/2$.

Pour $n=4$, chacun peut ouvrir deux boîtes. Les types de cycles favorables
sont
\[
1+1+1+1,\qquad 2+1+1,\qquad 2+2.
\]
Ils correspondent respectivement à $1$, $6$ et $3$ permutations. La
probabilité de réussite vaut donc
\[
\frac{1+6+3}{4!}=\frac{5}{12}.
\]

\subsection*{Comptage pour un entier pair $n$}

La stratégie échoue exactement lorsque la permutation contient un cycle de
longueur $k>n/2$. Pour une valeur fixée de $k$, le nombre de permutations
contenant un tel cycle est
\[
\binom nk (k-1)!(n-k)!=\frac{n!}{k}.
\]
En effet, on choisit les $k$ éléments du cycle, on les ordonne cycliquement de
$(k-1)!$ façons, puis on permute librement les $n-k$ éléments restants.

Une permutation ne peut pas contenir deux cycles de longueurs strictement
supérieures à $n/2$. Les événements comptés pour les différentes valeurs de
$k$ sont donc incompatibles. Pour une permutation uniforme,
\[
\Pp(\text{échec})
=\sum_{k=n/2+1}^{n}\frac1k.
\]
En notant $H_m=\sum_{k=1}^m 1/k$ le $m$-ième nombre harmonique, on obtient
\[
\Pp(\text{réussite})
=1-\sum_{k=n/2+1}^{n}\frac1k
=1-H_n+H_{n/2}.
\]

Pour $n=100$,
\[
\Pp(\text{réussite})
=1-\sum_{k=51}^{100}\frac1k
=1-H_{100}+H_{50}
\approx 0{,}3118278.
\]

\begin{resultbox}
La stratégie des cycles donne aux cent prisonniers une probabilité de réussite
collective d'environ
\[
\boxed{31{,}18\,\%},
\]
contre $2^{-100}$ seulement avec des choix aléatoires indépendants.
\end{resultbox}

Enfin, comme $H_n-H_{n/2}\to\log 2$, cette probabilité tend, lorsque $n$ pair
tend vers l'infini, vers
\[
1-\log 2\approx 30{,}69\,\%.
\]

\end{document}
